\chapter{Processing Pipeline and Identity Learning}

\section{Introduction}

This chapter explains how an uploaded video becomes a set of tracks, identities and analytics. It first describes the end-to-end flow (upload, process, retrieve), then each processing stage, the state models that describe progress, and the configuration parameters that control behaviour.

\section{End-to-End Flow}

\subsection{Upload}

The user posts a video to \code{POST /videos}. The backend saves the file through \code{storage_service.save_upload}, extracts its metadata with \code{extract_video_metadata}, and creates a \code{Video} row with status \code{uploaded}.

\subsection{Process}

A request to \code{POST /videos/{id}/process} first validates that there is no active job for the video. It then creates a \code{ProcessingJob} with status \code{QUEUED}, and either enqueues the work (Celery) or runs it synchronously as a background task. The background task performs:

\begin{enumerate}
  \item frame extraction according to the \code{PipelineConfig};
  \item object detection, using the model specified by \code{model_path};
  \item tracking, using the \code{mock}, \code{bytetrack} or \code{tapnext} tracker;
  \item re-identification embeddings;
  \item identity consolidation;
  \item analytics computation (\code{compute_video_analytics}).
\end{enumerate}

When processing finishes, the \code{Video} status becomes \code{processed}, and the \code{IdentityRecord}, \code{Track}, \code{DetectionRecord}, \code{FrameTimingRecord} and \code{CameraMotionRecord} rows are created.

\subsection{Retrieve}

The GET endpoints return the persisted tracks, identities, analytics, frame timings, motion estimates, ReID events and detections. Retrieval never triggers recomputation.

\begin{figure}[htbp]
\centering
\begin{tikzpicture}[
  node distance=6mm,
  st/.style={draw,rounded corners=2pt,align=center,font=\footnotesize,minimum height=9mm,text width=1.9cm},
  arr/.style={-{Stealth[length=2mm]},thick}]
\node[st] (a) {Upload video};
\node[st,right=of a] (b) {Frame extraction};
\node[st,right=of b] (c) {Object detection};
\node[st,right=of c] (d) {Tracking};
\node[st,below=8mm of d] (e) {Re-ID embeddings};
\node[st,left=of e] (f) {Identity consolidation};
\node[st,left=of f] (g) {Analytics};
\node[st,left=of g] (h) {Persist \& serve via API};
\draw[arr] (a) -- (b);\draw[arr] (b) -- (c);\draw[arr] (c) -- (d);
\draw[arr] (d) -- (e);\draw[arr] (e) -- (f);\draw[arr] (f) -- (g);\draw[arr] (g) -- (h);
\end{tikzpicture}
\caption{Processing pipeline from upload to persisted results.}
\label{fig:pipeline}
\end{figure}

\section{Status Models}

The progress of a video and of a processing job is described by the shared TypeScript types in \code{packages/types/src/common.ts}:

\begin{itemize}
  \item \code{JobStatus}: \code{uploaded} $\rightarrow$ \code{processing} $\rightarrow$ \code{processed}, or \code{error};
  \item \code{ProcessingStage}: \code{frame_extraction} $\rightarrow$ \code{detection} $\rightarrow$ \code{tracking} $\rightarrow$ \code{reid} $\rightarrow$ \code{analytics};
  \item \code{ProcessingStatus}: \code{QUEUED} $\rightarrow$ \code{RUNNING} $\rightarrow$ \code{COMPLETED}, or \code{FAILED}.
\end{itemize}

\begin{figure}[htbp]
\centering
\begin{tikzpicture}[
  node distance=14mm,
  s/.style={draw,rounded corners=6pt,align=center,font=\footnotesize,minimum height=8mm,minimum width=2.2cm},
  arr/.style={-{Stealth[length=2mm]},thick}]
\node[s] (q) {QUEUED};
\node[s,right=of q] (r) {RUNNING};
\node[s,right=of r,yshift=6mm] (c) {COMPLETED};
\node[s,right=of r,yshift=-6mm] (f) {FAILED};
\draw[arr] (q) -- (r);\draw[arr] (r) -- (c);\draw[arr] (r) -- (f);
\end{tikzpicture}
\caption{Life cycle of a \texttt{ProcessingJob} (\texttt{ProcessingStatus}).}
\label{fig:jobstate}
\end{figure}

\section{Processing Stages}

\subsection{Frame Extraction}

Frames are extracted from the uploaded video according to the \code{PipelineConfig} of the job. The extracted frames are stored through the \code{StorageService} in MinIO/S3, where the video's width, height, frame rate and frame count were recorded at upload time. \add{Sub-sampling the frame rate reduces cost on long videos but increases the displacement of fruit between sampled frames, so it directly affects tracker difficulty.}

\subsection{Object Detection}

Fruit are detected in each frame by the model specified by \code{model_path}. Detections below the configured confidence threshold are discarded, and the remainder are persisted as \code{DetectionRecord} rows with their bounding box and confidence.

\subsection{Tracking}

The tracker links detections across frames into tracks. The tracker type is selected from \code{mock}, \code{bytetrack} and \code{tapnext}; the tracker match threshold and the number of buffer frames (how long an unmatched track is kept alive) are configurable. \add{The \code{mock} tracker provides a deterministic back-end for development and tests, ByteTrack is a motion-based tracking-by-detection baseline, and TAPNext follows points through the video (see Chapter~2).}

\subsection{Re-Identification Embeddings}

For each tracked detection the pipeline computes a ReID embedding. The embeddings are used to compare observations with the prototypes of existing identities, using the ReID threshold to decide whether a match is acceptable.

\subsection{Identity Consolidation}

Identity consolidation turns tracks and embeddings into persistent identities. For each identity the pipeline stores a prototype embedding, a quality score, the observation count and an \code{is_confident} flag; every sighting is recorded as an \code{IdentityObservation}. When an identity is lost and later reappears, a \code{ReidEvent} stores the lost frame and the reappearance frame.

\add{A common formulation for maintaining the prototype $\mathbf{p}_k$ of identity $k$ is an exponential moving average of the embeddings assigned to it, controlled by the memory update rate $\alpha$:}

\begin{equation}
\colorbox{\phcol}{$\displaystyle \mathbf{p}_k \leftarrow \frac{(1-\alpha)\,\mathbf{p}_k + \alpha\,\mathbf{e}_t}{\left\lVert (1-\alpha)\,\mathbf{p}_k + \alpha\,\mathbf{e}_t\right\rVert_2}$}
\end{equation}

\add{A detection with embedding $\mathbf{e}_t$ is re-associated with identity $k$ when the cosine similarity between $\mathbf{e}_t$ and $\mathbf{p}_k$ exceeds the ReID threshold $\tau_{\mathrm{reid}}$. This equation illustrates the role of the memory update rate; the exact update rule should be confirmed against the implementation before submission.}

\subsection{Analytics}

The analytics stage runs \code{compute_video_analytics}, which uses the \code{AnalyticsEngine} to compute IDF1, HOTA and the number of ID switches. The aggregated analytics are served by \code{GET /videos/{id}/analytics}.

\section{Stage Inputs and Outputs}

\begin{table}[htbp]
\centering\small
\begin{tabular}{@{}>{\raggedright\arraybackslash}p{2.9cm}>{\raggedright\arraybackslash}p{3.7cm}>{\raggedright\arraybackslash}p{3.7cm}>{\raggedright\arraybackslash}p{3.1cm}@{}}
\toprule
\textbf{Stage} & \textbf{Input} & \textbf{Output} & \textbf{Persisted as}\\
\midrule
Frame extraction & Uploaded video, \code{PipelineConfig} & Frames & Frames in object storage\\
Detection & Frames, \code{model_path}, confidence threshold & Bounding boxes with confidence & \code{DetectionRecord}\\
Tracking & Detections, tracker type, match threshold, buffer frames & Tracks & \code{Track}\\
ReID & Tracked detections & Embeddings & (used by next stage)\\
Identity consolidation & Tracks, embeddings, ReID threshold, memory update rate & Identity prototypes, observations, ReID events & \code{IdentityRecord}, \code{IdentityObservation}, \code{ReidEvent}\\
Analytics & Tracks, identities & IDF1, HOTA, ID switches, timings, motion & \code{FrameTimingRecord}, \code{CameraMotionRecord}\\
\bottomrule
\end{tabular}
\caption[Stage inputs and outputs]{Inputs, outputs and persisted records of each pipeline stage. \add{The input and output columns are summarised from the pipeline description; they are not an exhaustive interface specification.}}
\label{tab:stages}
\end{table}

\section{Camera Motion}

The camera motion of each video is estimated and stored in a \code{CameraMotionRecord} with its motion type and whether compensation was applied. The camera-motion mode is one of the pipeline settings. \add{Compensating for ego-motion before association keeps motion predictions meaningful when the camera moves along an orchard row.}

\section{Identity-Space Visualisation}

The \code{EmbeddingProjection} service reduces identity prototypes to two dimensions with t-SNE or PCA, so that the interface can display them. The coordinates are served by \code{GET /videos/{id}/identities/projection}. \add{Identities of the same fruit should form tight clusters, and well-separated clusters suggest that the ReID threshold is set sensibly.}

\section{Identity Consolidation Procedure}

\add{The listing below gives a simplified, illustrative description of how tracks and embeddings can be consolidated into identities. It is provided to explain the idea and does not reproduce the source code of the implementation.}

\begin{lstlisting}[style=added,language=Python,mathescape=false]
for track in tracks_sorted_by_first_seen:
    e = mean_embedding(track)
    k, sim = nearest_identity(e, identities)   # cosine similarity
    if k is not None and sim >= reid_threshold:
        assign(track, identities[k])
        record_reid_event_if_gap(track, identities[k])
        update_prototype(identities[k], e, alpha)
    else:
        identities.append(new_identity(track, e))
for k in identities:
    k.quality = mean_quality(k.observations)
    k.is_confident = k.observation_count >= min_obs
\end{lstlisting}

\add{In this scheme a track that begins after another has ended is compared with all existing prototypes; if the similarity exceeds the ReID threshold it inherits the identity and a ReID event (lost frame, reappearance frame) is stored. Otherwise a new identity is created. Tracks that overlap in time must not be merged into the same identity, which acts as a simple constraint against merging two different fruit.}

\section{Design Considerations}

\begin{itemize}
  \item \add{\textbf{Threshold trade-off.} A low ReID threshold merges different fruit (undercounting); a high one fragments one fruit into several identities (overcounting). The threshold is therefore exposed as a setting and evaluated through experiments.}
  \item \add{\textbf{Prototype memory.} A small memory update rate makes prototypes stable but slow to adapt to lighting changes; a large one adapts quickly but can drift towards a neighbouring fruit.}
  \item \add{\textbf{Separation of stages.} Persisting detections, tracks, embeddings-derived identities and timings separately makes it possible to diagnose in which stage a counting error originated.}
  \item \add{\textbf{Tracker as a plug-in.} Because the mock, ByteTrack and TAPNext trackers are selected by a setting, a new tracker can be compared with the existing ones using the same video, detector and metrics.}
\end{itemize}

\section{Computational Considerations}

\add{Detection and embedding extraction dominate the processing time per frame; association and identity matching grow with the number of active tracks and identities. The per-frame timing records stored for each video allow these costs to be measured on the actual hardware, which is the basis of the runtime study in Chapter~8.}

\section{Pipeline Parameters}

\begin{table}[htbp]
\centering\small
\begin{tabular}{@{}>{\raggedright\arraybackslash}p{4.3cm}>{\raggedright\arraybackslash}p{9.5cm}@{}}
\toprule
\textbf{Parameter} & \textbf{Purpose}\\
\midrule
Confidence threshold & \add{Minimum detector confidence for a detection to be kept}\\
ReID threshold & \add{Minimum similarity for re-associating an observation with an existing identity}\\
Tracker type & Selects the tracker: \code{mock}, \code{bytetrack} or \code{tapnext}\\
Camera motion mode & \add{Selects how camera motion is estimated and compensated}\\
Memory update rate & \add{Controls how quickly an identity prototype adapts to new observations}\\
Tracker match threshold & \add{Minimum score for linking a detection to an existing track}\\
Buffer frames & \add{Number of frames an unmatched track is retained before it is considered lost}\\
\bottomrule
\end{tabular}
\caption{Key pipeline parameters (defaults stored in \texttt{AppSettings}).}
\label{tab:params}
\end{table}

\section{Failure Handling}

\add{A processing job can end in the \code{FAILED} state, and the corresponding video in the \code{error} state, for example when the video cannot be decoded or a stage raises an exception. Because every stage persists its output, the cause of a failure can be located from the \code{stage} recorded on the job. A failed video can be deleted with \code{DELETE /videos/{id}} and uploaded again after the problem is fixed.}

\section{Reproducibility}

\add{Each \code{ProcessingJob} stores the configuration it ran with in its \code{config} dictionary, and each experiment stores its tracker, scenario and number of frames together with the resulting metrics. Together with the exports in JSON or CSV, this makes it possible to repeat an evaluation and to attach the exact configuration to every number reported in Chapter~8.}

\section{Concurrency and Job Control}

A video can have only one active processing job: requesting processing while a job is active is rejected (\add{the API defines a \code{ConflictError} for this kind of situation}). The state of a job (\code{status}, \code{stage}, \code{progress}) can be polled through \code{GET /jobs/{id}}. When the work is enqueued through Celery it runs outside the API process; otherwise it runs as a background task of the API process.
